% interface=en output=pdftex

\environment mstyle

\starttext

\TitlePage{508}{0}{how to install\\\ConTeXt}{}{1}

\subject {Introduction}

Nowadays most \TEX\ users will use one of the many distributions that are
available for free or commercially. This means that much of the hard work of
installation is already taken care of. When installed properly, the user will
quite certainly have the binaries, hyphenation patterns, fonts and some macro
packages on the system.

This manual is therefore not targeted at installing \TEX, but focusses on how
to get \CONTEXT\ running within an existing distribution. When, after reading
this text, you still cannot get \CONTEXT\ working properly, we advice you to
contact the maintainer of your distribution, or to send your questions to the
\CONTEXT\ mailing list.

\subject{Unpacking the archives}

The \CONTEXT\ distribution consists of several zipped archives. You need to
unpack these to the appropriate directory of you file system. The archives
are zipped using the free \type {zip} program, and can be unzipped using its
counterpart. \footnote {Shareware programs like WinZip for MSWindows also
work well.} Just type:

\starttyping
unzip -a <archive>.zip
\stoptyping

Everything you need to get started can be found in the main \CONTEXT\
archive:

\starttabulate[|Tl|l|]
\NC cont-tmf.zip \NC the \CONTEXT\ sources and programs \NC \NR
\stoptabulate

There are some more archives you may want to install, like:

\starttabulate[|Tl|l|]
\NC cont-img.zip \NC a couple of extra resources              \NC \NR
\NC cont-exa.zip \NC the \EXAMPLE\ framework tools and styles \NC \NR
\stoptabulate

There are two more somewhat redundant archives. These can be useful for users
of packages other than \CONTEXT. If you are using the files from \type
{cont-tmf.zip}, it is not necessary to include the files from these archives.

\starttabulate[|Tl|l|]
\NC cont-ppc.zip \NC the \PPCHTEX\ only files                 \NC \NR
\NC cont-mpd.zip \NC the \METAPOST\ to \PDF\ conversion tools \NC \NR
\stoptabulate

To prevent duplication in files, we strongly advise to obey the path as coded
in the archives. The files in these archives will be unpacked into the
appropriate directories of the official \type {texmf} tree. Some files are
stored in the directories for generic \TEX\ files or the \LATEX\ package,
but those are not important for \CONTEXT\ users.

\starttabulate[|Tl|l|]
\NC texmf/tex/context/config  \NC some configuration files           \NC \NR
\NC texmf/tex/context/base    \NC all \CONTEXT\ core files, modules  \NC \NR
\NC texmf/tex/context/extra   \NC some goodies                       \NC \NR
\NC texmf/tex/context/sample  \NC a couple of sample files           \NC \NR
\NC texmf/tex/context/user    \NC user specific files                \NC \NR
\NC texmf/tex/context/third   \NC third party contributions          \NC \NR
\NC texmf/metapost/context    \NC the \METAPOST\ modules             \NC \NR
\NC texmf/context/config      \NC some configuratiom files           \NC \NR
\NC texmf/context/data        \NC data filed used by scripts         \NC \NR
\NC texmf/scripts/context     \NC the \PERL\ and \RUBY\ scripts      \NC \NR
\NC texmf/doc/context/base    \NC \CONTEXT\ documentation            \NC \NR
\stoptabulate

Users can best not put files in the \type {base}, \type {extra} and \type
{sample} directories. That way they can conveniently be removed and
reinstalled. The files in \type {third} are to be organized by author.
\footnote {You need to create this (yet empty) directory yourself.}

The way \PERL\ and \RUBY\ scripts are started differs per operating
system. For instance, on \UNIX, the \PERL\ scripts should be
installed without the \type {.pl} extension, because these scripts
and possibly other programs rely on these names. They should be
moved to the search path for binaries and scripts. On \MSWINDOWS\
running scripts is not supported by the operating system, although
in  \MSWINDOWS\ 2000 and later you can confortably associate
extensions (suffixes) with programs.  Most script environment
installers will set this association for you.  In the binary tree of
the \TEX\ distributions you can find stubs like \type {irun},
written by Fabrice Popineau. When copied to \type
{<scriptname>.exe}, this program launches the script with the same
name.

The initialization files for \TEXEXEC\ go into \type {texmf/context}
(\TETEX) or into the same directory as the binaries. When setting up
\TETEX, make sure you enable generation of the format files, by
uncommenting the lines that specify the \CONTEXT\ formats. You can
edit the \TETEX\ (and \FPTEX) configuration file using:

\starttyping
fmtutil
\stoptyping

In \TETEX\ you should pass the switch \type {--edit}. When you pass \type
{--all} you get all formats. When in need for patterns other than the default
ones, \type {texconfig} can be used to enable more hyphenation patterns in
the file \type {cont-usr.tex}. You can also edit this file directly.

Another way to deal with scripts is to use \TEXMFSTART\

\starttyping
texmfstart --make --windows --indirect texexec.pl
texmfstart --make --unix    --indirect texexec.pl
\stoptyping

The resulting files (\type {texexec} and/or \type {texexec.bat} can be moved
to your binary path. The advantage of this approach is that it is less sensitive
to changes in the formal directory structure, since it will be maintained in a
downward compatible way. You need to do the same for other tools, like \TEXUTIL.

\CONTEXT\ comes with a font that contains navigational symbols. These fonts
go to the corresponding places in the fonts tree, in our case:

\starttabulate[|Tl|l|]
\NC texmf/fonts/tfm/hoekwater/context   \NC the files with suffix \type{tfm} \NC \NR
\NC texmf/fonts/afm/hoekwater/context   \NC the files with suffix \type{afm} \NC \NR
\NC texmf/fonts/type1/hoekwater/context \NC the files with suffix \type{pfb} \NC \NR
\stoptabulate

There are also \CONTEXT\ specific encoding|/|map files. These are stored
in the following paths:

\starttabulate[|Tl|l|]
\NC texmf/fonts/map \NC \CONTEXT\ fontmap file \NC \NR
\NC texmf/fonts/enc \NC \CONTEXT\ font encoding file \NC \NR
\stoptabulate

In older distributions those files had to go under \type {texmf/dvips/config}.
Unfortunately \TEX\ distributions are not upward compatible in this respect.

Don't forget to add a reference to this file \type {context.map} to the \type
{pdftex.cfg} file that resides in the \PDFTEX\ configuration directory.

\starttyping
map +context.map
\stoptyping

Instead you can configure \type {cont-sys.tex} (being your personal copy of
\type {cont-sys.rme}) to let \CONTEXT\ load map files on demand.

\subject{Setting up \TEXEXEC}

\TEXEXEC\ is the command line interface to \CONTEXT. There is nothing wrong
with running \CONTEXT\ in the traditional way, like

\starttyping
pdfetex  &cont-en filename
\stoptyping

or for a \WEBC\ bases \TEX

\starttyping
pdfetex  &cont-en --progname=context filename
\stoptyping

but, and this will be more clear when we provide more options, the next call
is more convenient:

\starttyping
texexec  filename
\stoptyping

This not only hides the implementation specific switches, but also removes
the need for specifying the format.

\TEXEXEC\ is written in \PERL, a scripting language that is available on most
leading software platforms. In order to operate well, we need to set up
\TEXEXEC. Of course you must have \PERL\ running on your system. On \UNIX\
systems, this is always the case, but for \MSWINDOWS\ you have to install
\PERL\ yourself. \footnote {A good choice is the Active State distribution
that you can fetch from \type {www.activestate.com}}.

First you have to move \TEXEXEC\ and its relative \TEXUTIL\ to a location in
the binaries path. When issuing the command \type {texexec} you should get
some response. Even better, when saying:

\starttyping
texexec  --verbose
\stoptyping

you should get some information on how \TEXEXEC\ is set up. When generating
formats and processing files fail, you need to set up the initialization
file \type {texexec.ini}. This file comes disguised as \type {texexec.rme},
so when not present, you need to copy this file. The initialization file
should be present in the \type {config} path, or in the same path as the
script. Next you need to edit this file. The simplest way is to comment
and uncomment one of the following lines:

\starttyping
set  TeXShell  to  tetex
%set TeXShell  to  fptex
%set TeXShell  to  miktex
%set TeXShell  to  private
\stoptyping

When this is done, you should check to what extent the rest of the variables
in this file match the local settings. We hope that the names of the
variables used are clear.

When it is not set up properly, \TEXEXEC\ tries very hard to locate the files
it needs. Normally \TEXEXEC\ should start up rather quickly. When you are
under the impression that you are waiting too long, there is probably an
error in the setup.

\subject{Using \TETEX}

When you are using te\TEX\ or derived distributions, you can usually stick to
the regular updates, unless you want to use the latest version of \CONTEXT.
In any event, you should make sure that only one copy is present on your
system, because otherwise files can get mixed, due to the often aggressive
file searching algorithms. If you want to update anyway, you can unzip \type
{cont-tmf.zip} from within the \type{texmf} directory and regenerate the
format files.

The \PPCHTEX\ only archive is for users who maintain their own files and only
want to install this package. The \TEXUTIL\ archive is for those who want to
use this script, but don't want to install \CONTEXT. Neither archive is
needed when you install \CONTEXT\ from \type {cont-tmf.zip}.

\subject{Using \FPTEX}

The first step in installing \CONTEXT\ under \FPTEX, is to unzip the file
\type {cont-tmf.zip} within the \type {texmf} directory. Afterwards the \type
{texexec} binaries and \PERL\ script should be copied to the directories that
contain the other \TEX\ binaries. The \TEX\ binaries path, that should also
be part of your \type {PATH} variable, often looks like:

\starttyping
../tex/bin/win32
\stoptyping

When no file \type {texexec.exe} is found, you have to unzip \type
{texexec.zip} and copy \type {runperl.exe} to \type {texexec.exe}. Don't
forget to update the \type {ls-R} file database by running \type {mktexlsr}.

Next you need to locate the file \type {texmf/web2c/fmtutil.cnf}. In this
file, which contains documentation, you need to activate the \CONTEXT\
formats. Now you can run:

\starttyping
texexec --make
\stoptyping

\subject{Using \MIKTEX}

This section is provided by Grzegorz Sapijaszko and Ed L.~Cashin and concerns
the installation of \CONTEXT\ under \MIKTEX. Installation in the \MIKTEX\
environment isn't much different from the \TETEX\ installation. Nevertheless,
you should take a few steps to achieve good results. The first one is copying
a \type {texexec.rme} file to \type {texexec.ini} and uncommenting the lines
for \MIKTEX\ as follows:

\starttyping
%set TeXShell  to  tetex
%set TeXShell  to  fptex
set TeXShell  to  miktex
%set TeXShell  to  private
\stoptyping

In next step you should add a string \type {\context\perltk} to your
environment variable \type {PATH} (in autoexec.bat under \WINNX, or in Control
Panel in \WINNT), for example:

\starttyping
PATH=c:\miktex\context\perltk;
\stoptyping

If you are using languages other than Dutch, German or English, you should
uncomment the lines in the \type {cont-usr.tex} file for the hyphenation
patterns of the languages you need. After that you can refresh the filename
databases, for example, by using \quotation {Refresh Filename Databases} from
the \quotation{Start \vl\ Programs \vl\ MikTeX \vl\ Maintenance} menu or
\quotation {\MIKTEX\ \vl\ Options} menu, depending on the version you run.
You should also generate a format file. This is described in the next
section. Notice that you should have \PERL\ installed on your system.

After generating the formats you should copy the format file \type
{cont-xx.fmt} from \type {\context\perltk} to the directory where \MIKTEX\
is storing formats (usually \type {\localtexmf\miktex\fmt}). Another way is
to add a \type {\context\perltk} directory to \type {miktex.ini} file:

\starttyping
[MiKTeX]
...
...
;; Where MiKTeX searches for .fmt files.
FMTPath=.;%R\miktex\fmt;C:\localtexmf\context\perltk//
\stoptyping

As a convenience, you can copy the \type {runperl.exe} file from the \type
{cont-wrk.zip} archive to \type {texexec.exe}. You should make sure sure that
those binaries are in the same directory as the \PERL\ scripts.

\subject{Generating formats}

From its name you can deduce that \CONTEXT\ is written in the typographic
language \TEX. \CONTEXT\ is parameter driven, which means that users change
its behaviour by setting variables and changing keys. \CONTEXT\ comes with a
multi||lingual interface. Currently there are three such interfaces: Dutch,
English and German.

Users who want complete control can edit the file \type {cont-usr.tex} and
generate a format using the main file \type {context.tex}. Users who want an
Dutch, English or German format, can stick to the files named \type
{cont-nl}, \type {cont-en}, and \type {cont-de}. Again, by editing the file
\type {cont-usr.tex}, you can influence the outcome.

In the early years of \TEX, generating a format was common practice and users
were pretty well aware of format files, hyphenation patterns and fonts.
Nowadays, disributions take care of the more complicated issues, so users can
comfortably skip many nasty installation steps. To make life even more
comfortable, \CONTEXT\ comes with \TEXEXEC, a command line interface to
\TEX. When properly set up, this \PERL\ script can save you much time.

For instance, generating the three formats mentioned is accomplished by:

\starttyping
texexec  --make  en de nl
\stoptyping

When \TEXEXEC\ is set up properly, this command should work. Before you read
on, you should try to generate at least the English and Dutch format. The
Dutch format is needed by some of the postprocessing features built into
\TEXEXEC.

\starttyping
texexec  --make  en nl
\stoptyping

Additionally, you need to make the \METAFUN\ format. This is an extension to
the standard \METAPOST\ format which offers not only more features, but also
interfaces quite well to \CONTEXT.

\starttyping
texexec  --make metafun
\stoptyping

The formats associated with the interfaces default to the language of the
interface. This is all right for Dutch, English or German users, but Polish
and Czech users are worse off. For them a format file that defaults to their
own language makes more sense. Poles will like to say:

\starttyping
texexec  --make  --language=pl  --bodyfont=plr  en
\stoptyping

while Czech people will go for:

\starttyping
texexec  --make  --language=cz  --bodyfont=csr  en
\stoptyping

Both produce a format called \type {cont-en} with an English interface, but
the first one defaults to Polish hyphenation patterns and fonts, and the
second one to Czech ones. If wanted, you may pass a comma separated list of
languages.

\starttyping
texexec  --make  --language=pl,it,uk  --bodyfont=plr  en
\stoptyping

or, to generate a english interface format with Czech and Slovak patterns and
Czech||Slovak Computer Modern Roman fonts:

\starttyping
texexec  --make  --language=cz,sk,en  --bodyfont=csr  en
\stoptyping

At \PRAGMA\ we simply load all the languages:

\starttyping
texexec --alone --make --all
\stoptyping

Unfortunately, the hyphenation patterns are of hard coded in a format file
and cannot be loaded at run time. When patterns are needed other than the
ones loaded by default, you can consider adapting the file \type
{cont-usr.tex}. This file is loaded at format generation time. When for
instance Italian patterns are to be used, given that these are available
either in the file \type {lang-it.pat}, or in a file onto which this filename
is mapped, you should uncomment the line:

\starttyping
\installlanguage [\s!it] [\c!state=\v!start] % italian
\stoptyping

The strange looking \type{\s!} and \type{\c!} things are needed in order to
support multiple interfaces. Don't touch these!

When using \WEBC, in \type {texmf.cnf} some \CONTEXT\ specific memory
settings take place. When directly generating a format |<|i.e.\ when you're
not using the \TETEX\ initialization script or \TEXEXEC|>| you should supply
the program name: \type {-progname=context}

Make sure you read the manual to \TEXEXEC. Apart from the normal processing
of files, there are quite a few useful options: mode dependant processing,
output selection, generating booklets, typesetting contact sheets of figures,
manipulating \PDF\ files, and more.

\subject{Changing defaults}

The somewhat more run||time specific settings, like certain special drivers,
can be added to \type {cont-sys.tex}. This file is loaded at run time. For
instance, this file can can contain the line:

\starttyping
\setupoutput[pdftex]
\stoptyping

This commands tells \CONTEXT\ to produce \PDF\ output by default. For Y\&Y
and Acrobat support, you just say:

\starttyping
\setupoutput[dviwindo,acrobat]
\stoptyping

Of course you can also load location specific layout settings in this file.
The next few lines tell \CONTEXT\ to default to the \CONTEXT\ navigational
symbols, instead of the ones composed from other glyphs.

\starttyping
\usesymbols [nav]
\setupsymbolset [navigation 1]
\stoptyping

At \PRAGMA\ we want to process \METAPOST\ files at run||time, so there we
also have entries like:

\starttyping
\runMPgraphicstrue
\runMPTEXgraphicstrue
\useMETAFUNformattrue
\stoptyping

The first setting leads to less intermediate files. The next two settings let
\CONTEXT\ process the graphics real||time. For that you need to enable \type
{write18} in the \TEX\ configuration file (for \WEBC\ this is \type
{texmf.cnf}). The last setting only makes sense if you have also generated a
\METAFUN\ format. If you want to know more about \METAFUN, take a look at its
manual.

The verbatim environments provide pretty printing. When you want even more
fancy verbatim, for instance with in||between switching of a language
interpreter, you should enable this feature. If you get no idea what we're
talking about here, leave this variable untouched since tricky code is
involved.

\starttyping
\newprettytrue
\stoptyping

When \CONTEXT\ cannot determine the dimensions of an external figure, and no
\type {texutil.tuf} file is present, you can let \CONTEXT\ call \TEXUTIL\
directly. If you use \PDFTEX, you can leave this switch off.

\starttyping
\runutilityfiletrue
\stoptyping

More than one instance of \TEX\ a single path, can lead to clashes in
temporary files. The next switch enables a filename security feature:

\starttyping
\protectbufferstrue
\stoptyping

For the moment, we use these low level boolean switches instead of more
readable commands.

An important section in this file are the font defaults. You can set a
default font encoding as well as enable map file loading for \PDFTEX. Some of
these features are still experimental.

\subject{The \WEBC\ configuration}

Although not stricktly needed, \CONTEXT\ will operate more smoothly when
these features are set in the file {texmf.cnf}.

\starttyping
openout_any = a
shell_escape = t
allow_multiple_suffixes = f
\stoptyping

The first line permits \CONTEXT\ to open parent paths that can hold common
styles. The second line enables enables users to embed \METAPOST\ code in
their document. As a result, \METAPOST\ is called automatically. \footnote
{You can disable this feature by running \TEXEXEC\ with the switch \type
{--automp}.} The last line makes sure that when opening files like \type
{somefile.tuo}, \TEX\ will not try to open \type {somefile.tuo.tex} first.

When you embed \TEX\ code in a \METAPOST\ definition, using \typ {btex ...
etex}, the next line will use \TEXEXEC\ to process that fragment in the case
of using \METAPOST\ directly.

\starttyping
MPXCOMMAND=texexec --mptex
\stoptyping

However, the latest distribitions don't need that kind of tweaking any more
since \TEXEXEC\ sort things out itself.

\subject{Processing files}

You can test you installation with the following file.

\starttyping
\starttext
  \framed {Let's see if it works.}
\stoptext
\stoptyping

If \CONTEXT\ is set up properly, then the command

\starttyping
texexec  filename
\stoptyping

will produce a file \type {filename.dvi}. Unless a file was already processed
before, you will notice that \TEXEXEC\ processes the file at least two times.
During a \TEX\ run, \CONTEXT\ saves information in the files \type
{filename.tui}: cross references, entries to the table of contents, data
needed for optimization, etc. If the run is succesful, this file is
converted to a file called \type {filename.tuo} and used in the next run.
\TEXEXEC\ will reprocess the file until the \type {tuo} file is unchanged.

We {\em strongly recommend} to use \PDFETEX: Peter Breitenlohners \ETEX\
permits \CONTEXT\ to run more efficient, while \THANH's \PDFTEX\ provides
\PDF\ output and \PDFETEX\ combined both in one program. For those unfamiliar
with \ETEX: this is an extension of \TEX\ that is not only more efficient and
in some aspects faster, but that also offers a few features that come in
handy when writing complicated code. For the moment, \CONTEXT\ adapts its
behaviour to the kind of \TEX\ that is used, but future versions may rely on
\ETEX\ or other descendants completely.

By default, \CONTEXT\ generates \DVI\ output for \DVIPS, unless the output is
specified otherwise. We already mentioned the \type {\setupoutput} command. A
second way of achieving this is:

\starttyping
texexec  --pdf  filename
\stoptyping

And yet another way is adding a comment line in the document source, like:

\starttyping
% interface=en  output=pdftex  translate=cp1250pl
\stoptyping

or

\starttyping
% interface=en  output=pdftex  translate=cp1250cz
\stoptyping

Now we can omit the \type {--pdf} switch when we launch \TEXEXEC. Normally
\TEXEXEC\ is able to sort out the interface itself, but in case of troubles,
you can set some defaults in the file \type {texexec.ini}. The \type
{translate} key is only needed when you use the reencoding||on||the||fly
feature of \WEBC.

The distributions has a few demo files. You may try:

\starttyping
texexec --pdf demo-tex.tex
texexec --pdf demo-mps.tex
texexec --pdf demo-xml.tex
\stoptyping

% todo: xetex

\subject{Minimal \CONTEXT\ trees}

From the \PRAGMA\ website you can download a minimal \CONTEXT\
distribution. This is just a subset of \TEXLIVE. On \MSWINDOWS, this
works as follows:

\starttyping
mkdir c:\context
cd /d c:\context
unzip mswintex.zip
setuptex
texexec --make --alone --all
\stoptyping

You can now (whenever needed) initialize the tree with:

\starttyping
c:\context\setuptex.bat c:\context
\stoptyping

Normally you will do this when a shell (cmd) is opened. There is no
need to do additional configurations.

On \UNIX\ you can do something:

\starttyping
mkdir /home/myself/context
cd /home/myself/context
unzip mslinux.zip
. setuptex
texexec --make --alone --all
\stoptyping

Here you can initialize the tree with:

\starttyping
. /home/myself/context/setuptex /home/myself/context
\stoptyping

When you don't provide a path, this script will try to figure out things
itself. Since the tree is completely isolated from other installations,
you can comfortably use multiple trees alongside.

\subject{Subscribing to the list}

There are two dedicated mailing lists hosted by the \NTG:

\starttabulate[|l|l|]
\NC the \CONTEXT\ mailing list \NC \goto{\url[context]}[URL(context)] \NC \NR
\NC the \PPCHTEX\ mailing list \NC \goto{\url[ppchtex]}[URL(ppchtex)] \NC \NR
\stoptabulate

These lists are so called {\em mailman} lists. You can subscribe yourself to a
list using the web interface:

\starttyping
http://www.ntg.nl/mailman/listinfo/ntg-context
\stoptyping

More information can be found at:

\starttyping
http://www.pragma-ade.com
\stoptyping

Another source, rapidly growing in size and resources, is:

\starttyping
http://contextgarden.net
\stoptyping

\ColofonPage

\stoptext
